% ============================================================================
% Hypothesis H4: context-dependent dominance of climate and techno-economic
% uncertainty in the spatial site ranking.
% Pre-registered decision rule (OQ-050, H4), fixed 2026-10-13, before any
% real-country result was produced. Classification thresholds 0.40 and 0.60;
% strict cross-country rule; verdict per technology.
% Notation follows docs/METHODOLOGY.md and the H1 to H3 sections.
% Citation keys marked "verify" must be checked before submission.
% ============================================================================

\section{Hypothesis H4: Context-Dependent Dominance of Climate and
Techno-Economic Uncertainty}
\label{sec:h4}

\subsection{Purpose and Position in the Research Design}
\label{subsec:h4_purpose}

Hypothesis H4 asks whether the source of uncertainty that most changes the
spatial ranking of candidate sites, climate or techno-economic, is the same in
every country. It states that it is not: for a given technology, the climate
axis dominates the displacement of the ranking in some of the countries studied
and the techno-economic axis in others, so that no country-independent ordering
of the two sources holds.

H1 to H3 are evaluated for each country and technology. H4 is comparative across
countries and is evaluated for each technology. It answers the question of
generalisation (RQ5) and relies on the separability of the two uncertainty axes
(RQ4): each axis is introduced alone, with the other held at its nominal value,
and the displacement each produces in the site ranking is compared. H3 measured
the displacement produced by the two axes together; H4 attributes it.

The partitioning of uncertainty by source is established practice in climate
science, where the contributions of model, scenario and internal variability to
projection uncertainty are separated \cite{hawkinssutton2009}, and in
sensitivity analysis, where factor prioritisation identifies the inputs that
most affect an output \cite{saltelli2008,pianosi2016}. H4 applies the same
logic with one change: the output is not a model quantity such as the
levelised cost of energy, but the site selection that the ranking produces. The
design follows the Decision Making under Deep Uncertainty perspective, in which
climate futures are explored without probabilities
\cite{lempert2003,marchau2019}.

\subsection{Why the Ranking and Not the Cost Variance}
\label{subsec:h4_ranking}

A source of uncertainty can change the cost of every site substantially without
changing which sites are best. If an uncertain parameter transforms the
levelised cost of all sites by the same increasing function, for instance
$\mathrm{LCOE}'_i=a+b\,\mathrm{LCOE}_i$ with $b>0$, the order of the sites is
preserved:
\begin{equation}
    \mathrm{LCOE}_i<\mathrm{LCOE}_j\iff\mathrm{LCOE}'_i<\mathrm{LCOE}'_j .
    \label{eq:h4_rank_invariance}
\end{equation}
Capital cost and the discount rate act almost in this way, because they scale the
cost of every site in a country. In a decomposition of the variance of the
levelised cost, the techno-economic axis would therefore appear dominant
everywhere, almost by construction, while changing few siting decisions. Only
the parameters whose effect differs between sites, such as the cost per
kilometre of grid and road connection or the temperature coefficient of the
modules interacting with local temperature, move the ranking. A variance-based
attribution answers a different question from the one posed here, and is
reported only as secondary evidence. The primary statistic is defined on the
selection.

\subsection{Notation}
\label{subsec:h4_notation}

The notation of Sections~\ref{sec:h1} to~\ref{sec:h3} applies, for one country
and one technology. Both axis-specific rankings and the nominal ranking are
computed over the same set $\mathcal{I}$; $q_i$ is the installable capacity of
the central land scenario; $k=\lceil p_k N\rceil$ with $N=|\mathcal{I}_0|$; and
$m=k$. The nominal rank $r^{\mathrm{nom}}_i$ is the rank of cell $i$ within
$\mathcal{I}$ by $\mathrm{LCOE}_{i,f_0}$, as in H3.

Two axis-isolated scores are defined from the relative regret $r_{i,f}$ of
Equation~\ref{eq:h1_relative_regret}.

Climate axis. Techno-economic parameters are held at their nominal values and
the climate varies over all members $m$ of the core window:
\begin{equation}
    Z^{\mathrm{clim}}_i=\max_{m}\; r_{i,(m,s_0)} .
    \label{eq:h4_clim_score}
\end{equation}

Techno-economic axis. The climate is held at the reference $m_0$ and the
parameters vary over the samples:
\begin{equation}
    Z^{\mathrm{tec}}_i=\mathrm{P}_{90,\,s\geq1}\!\left(r_{i,(m_0,s)}\right) .
    \label{eq:h4_tec_score}
\end{equation}

Each score keeps the operator that the primary robustness metric applies to its
axis: the worst case over the climate members, which carry no probabilities,
and the 90th percentile over the parameter samples, which form a sampled
probabilistic space. With both axes active, the two operators combined give the
primary metric $\mathrm{MR}_i$.

For each axis $a\in\{\mathrm{clim},\mathrm{tec}\}$, $S^{a}_k$ is the set of the
$k$ cells of $\mathcal{I}$ with the lowest $Z^{a}_i$. Ranks follow the
deterministic tie convention of the F7 phase record.

\subsection{Axis Displacement}
\label{subsec:h4_displacement}

The displacement produced by axis $a$ is the statistic of H3 applied to the
axis-isolated shortlist:
\begin{equation}
    D_a=\frac{\sum_{i\in S^{a}_k}q_i\,\mathbb{1}\!\left[r^{\mathrm{nom}}_i>k+m\right]}
    {\sum_{i\in S^{a}_k}q_i},
    \qquad a\in\{\mathrm{clim},\mathrm{tec}\}.
    \label{eq:h4_displacement}
\end{equation}
$D_a$ is the share of the capacity of the axis-isolated shortlist located in
cells that the nominal ranking does not reach even with the tolerance $m=k$.
Movements within the top $k+m$ are boundary fluctuation and are not counted.
$0\leq D_a\leq1$.

\subsection{Dominance Ratio and Country Profile}
\label{subsec:h4_profile}

The dominance ratio is the climate share of the combined axis displacement:
\begin{equation}
    \rho=\frac{D_{\mathrm{clim}}}{D_{\mathrm{clim}}+D_{\mathrm{tec}}},
    \qquad D_{\mathrm{sum}}=D_{\mathrm{clim}}+D_{\mathrm{tec}}>0 .
    \label{eq:h4_ratio}
\end{equation}
Each country and technology receives one profile:
\begin{equation}
    \mathrm{Profile}=
    \begin{cases}
        \text{No Displacement}, & D_{\mathrm{sum}}=0,\\
        \text{Climate-Dominated}, & \rho>0.60,\\
        \text{Techno-Economic-Dominated}, & \rho<0.40,\\
        \text{Mixed}, & 0.40\leq\rho\leq0.60 .
    \end{cases}
    \label{eq:h4_profile}
\end{equation}
The interval $[0.40,0.60]$ prevents a dominance label when the two axes displace
similar amounts of capacity; values exactly at 0.40 or 0.60 are Mixed. A country
whose inputs do not allow both displacements to be computed is reported as
unavailable and receives no profile.

The ratio is relative: it can be high when both displacements are small.
$D_{\mathrm{clim}}$, $D_{\mathrm{tec}}$ and $D_{\mathrm{sum}}$ are therefore
reported with every profile, and the profile is read together with them.

\subsection{Hypothesis and Decision Rule}
\label{subsec:h4_rule}

For a technology $\tau$, let $\mathcal{N}_\tau$ be the countries with a profile.
\begin{equation}
    \begin{aligned}
        \text{H4 is supported for }\tau \iff{} & \exists\, n_1\neq n_2\in\mathcal{N}_\tau:\\
        & \mathrm{Profile}(n_1)=\text{Climate-Dominated},\\
        & \mathrm{Profile}(n_2)=\text{Techno-Economic-Dominated}.
    \end{aligned}
    \label{eq:h4_rule}
\end{equation}
If fewer than two countries have a profile for technology $\tau$, H4 is not
testable for that technology, since a single national context cannot show a
difference between contexts. When at least two countries have a profile and no
climate-dominated and techno-economic-dominated pair exists, H4 is not
supported.
The rule is strict. Combinations of one dominance class with Mixed or No
Displacement countries do not support H4, because a Mixed profile does not show
a different dominant axis but the absence of one. The verdict is issued for
solar photovoltaics and for onshore wind separately; a combined statement across
the two technologies is descriptive.

The rule is deterministic and categorical: it is not a significance test and
does not estimate the probability that either axis will dominate in the future.
With three countries, a supported verdict shows that the dominant axis differs
in at least two of the national contexts studied; it does not establish a
pattern beyond them.

The thresholds 0.40 and 0.60 are decision thresholds of this thesis. The
literature on uncertainty partitioning and sensitivity analysis motivates the
attribution of displacement to sources but provides no universal boundary for
dominance. Their defence rests on pre-registration, on their reading (one axis
accounts for more than three-fifths of the combined displacement) and on the
reporting of the continuous ratio.

An outcome in which every country is Techno-Economic-Dominated is plausible,
because climate enters GeoFREA through annual-mean changes of the resource. It
would not support H4, and it would be a substantive finding: for siting up to
the core window, techno-economic uncertainty would outweigh climate change in
all the contexts studied.

\subsection{Complementary Evidence (No Verdict)}
\label{subsec:h4_complementary}

\begin{itemize}
    \item Interaction. The displacement of the joint shortlist, $D$ of H3, is
    reported beside $D_{\mathrm{clim}}$ and $D_{\mathrm{tec}}$. A joint
    displacement different from what either axis produces alone indicates that
    the two axes interact in the ranking.
    \item Variance shares. The decomposition of the variance of the relative
    regret over futures into a climate part and a techno-economic part, by the
    law of total variance, is reported as secondary evidence, with the caveat of
    Section~\ref{subsec:h4_ranking}.
    \item Scenario detail. $D_{\mathrm{clim}}$ recomputed with the climate members
    of each SSP separately, to show which scenario drives the climate axis.
    \item Distance to the boundaries. $|\rho-0.40|$ and $|\rho-0.60|$, which
    identify profiles close to a change of class.
\end{itemize}

\subsection{Sensitivity Analyses}
\label{subsec:h4_sensitivity}

The verdict is issued from the primary configuration only; a change of profile or
verdict in a variant is reported as a dependence on that choice.
\begin{enumerate}
    \item Without IPSL-CM6A-LR in the climate axis.
    \item The 2071--2100 window for the climate axis, in which the scenarios
    diverge more strongly.
    \item The 0.1-degree cell scale.
    \item Top-$k$ share $p_k\in\{0.05,0.10,0.20\}$ and tolerance
    $m\in\{k/2,\,2k\}$.
    \item Classification boundaries $\{0.30,0.70\}$ and $\{0.45,0.55\}$, reported
    as descriptive variants only.
\end{enumerate}

\subsection{Worked Example}
\label{subsec:h4_example}

Three countries are evaluated for solar photovoltaics. The values are
illustrative and are not GeoFREA results.

\begin{table}[htbp]
    \centering
    \caption{Illustrative axis displacements and profiles for H4.}
    \label{tab:h4_example}
    \small
    \begin{tabular}{lccccl}
        \toprule
        \textbf{Country} & $D_{\mathrm{clim}}$ & $D_{\mathrm{tec}}$ &
        $D_{\mathrm{sum}}$ & $\rho$ & \textbf{Profile} \\
        \midrule
        A & 0.12 & 0.38 & 0.50 & 0.24 & Techno-Economic-Dominated \\
        B & 0.42 & 0.08 & 0.50 & 0.84 & Climate-Dominated \\
        C & 0.25 & 0.25 & 0.50 & 0.50 & Mixed \\
        \bottomrule
    \end{tabular}
\end{table}

Country A is techno-economic-dominated and country B climate-dominated, so H4 is
supported for solar photovoltaics. In country A, getting the connection costs
right matters more than the climate for choosing where to build; in country B
the reverse holds. Neither statement implies that the other axis is negligible
in that country: the profile says which axis accounts for the larger share of
the combined displacement. If country B had $\rho=0.55$, the three profiles
would be techno-economic, Mixed and Mixed, and H4 would not be supported.

\subsection{Parameters}
\label{subsec:h4_parameters}

\begin{table}[htbp]
    \centering
    \caption{Pre-registered parameters of H4.}
    \label{tab:h4_parameters}
    \small
    \renewcommand{\arraystretch}{1.2}
    \begin{tabularx}{\textwidth}{
        @{}
        >{\raggedright\arraybackslash}p{2.6cm}
        >{\raggedright\arraybackslash}p{3.0cm}
        >{\raggedright\arraybackslash}X
        @{}
    }
        \toprule
        \textbf{Parameter} & \textbf{Value} & \textbf{Role} \\
        \midrule
        $Z^{\mathrm{clim}}_i$ & $\max_m r_{i,(m,s_0)}$ & Climate-only score;
        worst case over the core-window members. \\
        \addlinespace
        $Z^{\mathrm{tec}}_i$ & $\mathrm{P}_{90,s\geq1}\,r_{i,(m_0,s)}$ &
        Techno-economic-only score; 90th percentile over samples. \\
        \addlinespace
        $D_a$ & Statistic of H3 & Capacity of the axis shortlist beyond the
        nominal top $k+m$. \\
        \addlinespace
        $m$ & $k$ & Tolerance, as in H1b and H3. \\
        \addlinespace
        Boundaries & 0.40 and 0.60 & Techno-economic below 0.40, climate above
        0.60, Mixed between, inclusive. \\
        \addlinespace
        Verdict & Strict, per technology & At least one climate-dominated and
        one techno-economic-dominated country. \\
        \bottomrule
    \end{tabularx}
\end{table}

\subsection{Limitations}
\label{subsec:h4_limitations}

Isolation of the axes. Holding one axis at its nominal value measures each
source alone and does not capture their interaction, which is reported through
the joint displacement of H3 but not attributed.

Dependence on the nominal reference. Both displacements are measured against the
nominal ranking at $f_0$; a different reference would change the cells counted
as displaced.

Dependence on the uncertainty representation. The climate axis depends on the
six climate models, the three scenarios and the representation of climate
change through annual-mean change factors of the resource; the techno-economic
axis depends on the parameter ranges and their independent sampling. A narrower
climate representation tends to understate the climate axis.

Relative measure. A dominance profile can be assigned when both displacements are
small; the absolute values are reported to prevent over-interpretation.

Few countries. Three national contexts can show that the dominant axis is not
universal, but they cannot characterise how it varies with geography or explain
why; attributing a profile to a physical or economic mechanism requires evidence
beyond the classification.
